
% ════════════════════════════════════════
% 模型的评价（07.模型评价.tex）
% ════════════════════════════════════════
\section{模型的评价}

\subsection{模型的优点}

模型以线性规划为核心，决策变量全部连续，求解器可严格保证全局最优，不存在启发式算法陷入局部最优的问题。滚动调整模型用上境变量将分段惩罚费用精确线性化，不引入整数变量即完整刻画违约与超额的经济代价。逐日滚动与储电量传递机制把全年大规模优化分解为数百个日尺度小问题，在降低计算复杂度的同时保持全年储电量的连续性。该结构只依赖电价与不可控出力两列输入，替换序列即可复用，其稳健性已由波动电价情形与灵敏度分析验证。

\subsection{模型的缺点}

模型假定各 10 分钟间隔内功率恒定（假设1），分钟级剧烈波动造成的区间功率不平衡仍由紧急购电吸收，费用因此被低估；储能严格按计划执行的假定（假设5）同样忽略了储能对实时偏差的闭环补偿，实际紧急购电费应低于模型估计值。情景由历史残差的经验分布生成，仅取低、中、高三档共 9 个情景，对预测误差尾部刻画不足，连续阴雨等极端天气下的紧急购电费偏低。此外模型未计入电池循环寿命损耗与自放电，长周期运行下储能的实际经济性低于计算结果，容量配置与折旧成本之间的权衡尚未纳入优化。
